%% kapa: exam=2 round=37 year=2026 period=3 subject=compensation_law subject_ko=감정평가 및 보상법규 seq=2 number=2 points=30
%% source: Q-Net 감정평가사 기출문제 (cst00309) artlSeq=5268128 file=raw/1008/5268128/26년 제37회 감정평가사 2차 (3교시_감정평가 및 보상 법규).pdf
%% conversion: kordoc 4.24.51 -> tools/gen2.py (auto); review=proofread-gpt-terra; flags=none
\begin{longquestion}{【문제 2】}
국토교통부장관은 「부동산 가격공시에 관한 법률」(이하 ‘부동산공시법’ 이라 한다)에 따라 2025. 1. 24. X도 Y시에 소재한 이 사건 토지에 대해서는 그 가격을 m\textsuperscript{2} 당 700만원(이하 ‘이 사건 공시지가’라 한다)으로 결정하고 공시하였다. 甲은 2025. 2. 5. 이 사건 토지를 취득하였고, Y시장은 甲에게 2025. 9. 12. 이 사건 공시지가를 시가표준액으로 하여 산정한 재산세를 부과하였다(이하 ‘이 사건 과세처분’이라 한다). 한편, 2025. 7. 경 수사기관의 수사 결과 이 사건 공시지가의 조사ㆍ평가를 의뢰받은 A감정평가법인(이하 ‘A법인’이라 한다)의 소속 감정평가사로서 해당 업무를 수행한 B가 당시 이 사건 토지의 소유자였던 C로부터 금품을 받고 C의 요청에 근접한 값으로 이 사건 토지의 적정가격을 평가한 조사ㆍ평가보고서를 작성하였고, 국토교통부장관은 A법인이 제출한 이 조사ㆍ평가보고서에 기재된 내용대로 이 사건 공시지가 결정을 한 것으로 드러났다. 다음 물음에 답하시오. (30점) \par
\medskip\noindent 물음 1) 甲은 2025. 5. 7. 이 사건 공시지가 산정에 문제가 있음을 인지하고 감정평가사 D에게 의뢰하여 2025. 7. 1. 기준 이 사건 토지의 가액을 m\textsuperscript{2} 당 400만원으로 평가받은 바 있다. 먼저 甲이 부동산공시법 또는 「행정소송법」에 따라 이 사건 공시지가 결정을 직접 다툴 수 있는 방법을 설명하고, 甲이 2025. 11. 14. 이 사건 과세처분에 대하여 행정소송을 제기한 경우 이 사건 공시지가 결정의 하자를 주장할 수 있는지를 검토하시오. (20점) \par
\medskip\noindent 물음 2) A법인 및 B에게 적용될 수 있는 「감정평가 및 감정평가사에 관한 법률」상의 벌칙 및 행정상 제재(징계 포함)의 종류를 설명하시오. (10점) \par
\begin{framedtext}
참조 조문 \par
<지방세법> \par
제4조(부동산 등의 시가표준액) ① 이 법에서 적용하는 토지 및 주택에 대한 시가표준액은 「부동산 가격공시에 관한 법률」에 따라 공시된 가액(價額)으로 한다. 다만, 개별공시지가 또는 개별주택가격이 공시되지 아니한 경우에는 특별자치시장ㆍ특별자치도지사ㆍ시장ㆍ군수 또는 구청장(자치구의 구청장을 말한다. 이하 같다)이 같은 법에 따라 국토교통부장관이 제공한 토지가격비준표 또는 주택가격비준표를 사용하여 산정한 가액으로 하고, 공동주택가격이 공시되지 아니한 경우에는 대통령령으로 정하는 기준에 따라 특별자치시장ㆍ특별자치도지사ㆍ시장ㆍ군수 또는 구청장이 산정한 가액으로 한다. \par
② ∼ ④ (생 략) \par
제110조(과세표준) ① 토지ㆍ건축물ㆍ주택에 대한 재산세의 과세표준은 제4조제1항 및 제2항에 따른 시가표준액에 부동산 시장의 동향과 지방재정 여건 등을 고려하여 다음 각 호의 어느 하나에서 정한 범위에서 대통령령으로 정하는 공정시장가액비율을 곱하여 산정한 가액으로 한다. \par
1. 토지 및 건축물: 시가표준액의 100분의 50부터 100분의 90까지 \par
2. (생 략) \par
② ∼ ③ (생 략) \par
\end{framedtext}
\end{longquestion}
